\chapter{Observables and Spectral Decomposition}
\label{ch:observables}

Chapter~\ref{ch:born-rule} established probabilities for a fixed orthonormal
context, $p_i=\operatorname{tr}(\rho P_i)$ with $P_i=|e_i\rangle\langle e_i|$.
This chapter asks what an \emph{observable} is, and why Hermitian operators
are the right general class — not by declaration, but via an explicit
operational bridge. The route runs in two steps: first widen the rank-one
projectors of Chapter~\ref{ch:born-rule} to allow coarser, physically
realistic measurements with grouped outcomes, then ask what it means to
attach a numerical value to each outcome of such a measurement. That second
step is where Hermitian operators enter — as a consequence of the
attachment, not as a new ingredient assumed alongside it.

\section{Sharp Measurements and Coarse-Grained Outcomes}

\begin{definition}[Sharp measurement]
\label{def:sharp-measurement}
\index{Sharp measurement (definition)}
A \emph{sharp measurement} is a family of mutually orthogonal projectors
$\{P_a\}$ with $\sum_a P_a = I$ ($P_aP_b=\delta_{ab}P_a$), representing
mutually exclusive, jointly exhaustive outcomes.
\end{definition}

Chapter~\ref{ch:born-rule} proved $p_i=\operatorname{tr}(\rho P_i)$ only for
\emph{rank-one} $P_i$ (a full orthonormal basis). The extension to
Definition~\ref{def:sharp-measurement}'s coarser projectors costs nothing new:

\begin{proposition}[Coarse-grained Born rule]
\label{prop:coarse-born}
\index{Coarse-grained Born rule (proposition)}
For a sharp measurement $\{P_a\}$ where each $P_a$ is a sum of rank-one basis
projectors, $P_a = \sum_{i\in S_a} P_i$ for disjoint index sets $S_a$
partitioning $\{1,\dots,d\}$, $\operatorname{tr}(\rho P_a) = \sum_{i\in S_a}
p_i$.
\end{proposition}

\begin{proof}
Immediate from linearity of trace and Theorem~\ref{thm:density-born}.
\end{proof}

\section{The Operational Bridge}

\begin{definition}[Observable]
\label{def:observable}
\index{Observable (definition)}
Given a sharp measurement $\{P_a\}$ and a real value $a$ assigned to each
outcome, the corresponding \emph{observable} is
\begin{equation}
A = \sum_a a\, P_a.
\end{equation}
\end{definition}

Since each $P_a$ is Hermitian and $a\in\R$, $A$ is Hermitian. This is one
direction of the bridge: packaging a real-labeled sharp measurement produces
a Hermitian operator, by construction, not by postulate.

\begin{theorem}[Spectral theorem, cited]
\label{thm:spectral}
\index{Spectral theorem, cited (theorem)}
Every Hermitian operator $A$ on a finite-dimensional complex Hilbert space
admits a decomposition $A = \sum_a a P_a$ with real $a$ and mutually
orthogonal projectors $P_a$ summing to $I$.
\end{theorem}

This is a classical result of linear algebra, entirely independent of this
book's reconstruction, and is not re-derived here; it is cited exactly as
Chapters~\ref{ch:birkhoff} and~\ref{ch:local-tomography} cited
Birkhoff--von Neumann's antecedents and Araki's parameter-counting argument.
What Theorem~\ref{thm:spectral} supplies is the \emph{converse} of
Definition~\ref{def:observable}: every Hermitian $A$ unpacks into exactly the
kind of real-labeled sharp measurement the definition packages.

\begin{remark}
\label{rem:real-labels-assumption}
The bridge is therefore an equivalence, not a declaration: real-valued sharp
projective measurements and Hermitian operators are the same data, related by
Definition~\ref{def:observable} in one direction and
Theorem~\ref{thm:spectral} in the other. "Observables are Hermitian" is the
statement of this equivalence, not an independent postulate layered on top of
the state-space structure of Parts I--II. But the equivalence has a hidden
cost worth naming rather than absorbing silently: it holds only because
Definition~\ref{def:observable} already assumed outcomes are labeled by
\emph{real} numbers. Nothing established so far forces this — it is a
modeling choice (numerical outcomes are what get recorded, compared, and
averaged), consistent with the book's discipline of not letting assumptions
go unmarked. Complex or otherwise non-real labelings are mathematically
available and simply produce non-Hermitian (but still well-defined)
operators; the restriction to real labels is what makes $A$ Hermitian, not a
fact independently established about the physical world.
\end{remark}

\section{The Eigenvalue Rule as a Corollary}

\begin{proposition}[Eigenvalue rule]
\label{prop:eigenvalue-rule}
\index{Eigenvalue rule (proposition)}
If $A|\psi\rangle = a|\psi\rangle$ for normalized $|\psi\rangle$, then
$\Pr(a\mid\psi) = 1$.
\end{proposition}

\begin{proof}
$A|\psi\rangle=a|\psi\rangle$ places $|\psi\rangle$ in the $a$-eigenspace,
which is exactly the range of the spectral projector $P_a$
(Theorem~\ref{thm:spectral}); since $|\psi\rangle$ is normalized and lies
entirely in that eigenspace, $P_a|\psi\rangle = |\psi\rangle$. By
Theorem~\ref{thm:pure-state-born} (equivalently
Proposition~\ref{prop:coarse-born} applied to a pure state),
$\Pr(a\mid\psi) = \langle\psi|P_a|\psi\rangle = \langle\psi|\psi\rangle = 1$.
\end{proof}

No new probability postulate was needed to get here. "An eigenstate returns
its eigenvalue with certainty" — a rule ordinarily stated as its own
axiom — turns out to be a corollary of Chapter~\ref{ch:born-rule}'s
already-proved Born rule, once Definition~\ref{def:observable} pins down
what "the eigenvalue" means operationally: nothing more than which spectral
projector a given state happens to lie inside.

\begin{ledgerentry}[End of Chapter~\ref{ch:observables}]
\textbf{Established:} an operational equivalence between real-valued sharp
projective measurements and Hermitian operators
(Definition~\ref{def:observable}, Theorem~\ref{thm:spectral}); the
eigenvalue rule as a corollary of Chapter~\ref{ch:born-rule}'s Born rule, not
an independent postulate (Proposition~\ref{prop:eigenvalue-rule}).\\
\textbf{Not established:} why outcome labels must be real (flagged explicitly
as an assumption, not derived, Remark~\ref{rem:real-labels-assumption}); the
spectral theorem itself is cited from classical linear algebra, not proved
within this book's reconstruction; compatible observables, simultaneous
diagonalization, and commutators are untouched; the status of general
(non-Hermitian, non-normal) operators is not addressed. Continuous dynamics
— how observables and states evolve in time — is the subject of the next
chapter.
\end{ledgerentry}
